\section{引言}
\label{sec:introduction-zh}

后调性作曲常从明确的符号材料出发，例如音级集合、音程类轮廓、十二音序列及其变换、音区轨迹或具有特征性的节奏表层。计算机辅助作曲环境长期支持这类结构的操作，但近年的神经音乐系统更多在调性续写、流行歌曲结构、演奏 MIDI 或音频生成任务上进行评估 \citep{assayag1999openmusic,briot2020deep}。这些设置没有直接回答当代乐谱作曲中的一个实际问题：生成模型能否在保留具名后调性约束的同时给出可编辑草图，并说明约束在哪些位置得到满足或被违反？

这一问题需要不同的交互方式和评估词汇。音级集合理论提供正常序、原型和音程向量 \citep{forte1973structure,straus2016posttonal}；序列技法要求明确处理 $P/R/I/RI$ 变换和十二音聚合；节奏与姿态意图则应在乐谱时间中表达，而不是依赖音频帧。输出还必须能够以记谱形式检查。因此，仅生成流畅的标记序列并不充分；若序列不能转换为乐谱，或无法解释其与指定材料的关系，便难以服务于实际作曲。

数据合法性构成第二项挑战。大量 1945 年后的相关作品仍受版权保护，抓取此类语料难以公开和复现。本文改用确定性的程序化语料作为受控测试床。该语料不用于模仿具体作曲家，也不用于建立历史风格模型，而是提供已知的条件--目标关系，使研究者能够在不传播受保护乐谱的前提下测量约束遵循、失败模式和解码权衡。

本文研究一种条件前缀 Transformer 和神经符号解码器。前者生成量化的乐谱事件序列，后者利用不可微的后调性诊断指标对采样候选进行排序。神经模型本身采用标准架构，不提出新的注意力结构。本文贡献在于把因果符号生成与音级集合、序列、节奏、姿态、音区及记谱层面的检查连接起来，并同时返回 MusicXML 草图和机器可读的解释。

研究回答以下三个问题：
\begin{itemize}
  \item \textbf{RQ1：}在共享同一训练生成器和固定测试条件时，四候选符号选择相对于与其对齐的第一个候选，会如何改变显式约束指标？
  \item \textbf{RQ2：}从条件前缀中隐藏音级集合、序列、节奏或姿态字段，会如何影响模型对未改变的隐藏目标的遵循？
  \item \textbf{RQ3：}生成片段能否导出为结构有效、具有指定小节和声部数的 MusicXML，并同时提供可检查的分析报告？
\end{itemize}

本文主要贡献如下：
\begin{enumerate}
  \item 构建一个合法、可复现、包含 24,000 个片段的乐谱层面后调性程序化语料方案；
  \item 设计紧凑的条件前缀与事件表示，并训练可在本地运行的 Transformer，以生成 4--16 小节、2--8 声部的片段；
  \item 提出基于音级集合、序列、节奏、姿态和音域诊断的推断期候选排序目标；
  \item 建立带来源门控的评估与导出层，将汇总指标连接到 MusicXML 样例和逐片段 JSON 报告。
\end{enumerate}

本文结论有明确边界。量化结果针对合成约束测试床，而非风格真实性。项目已准备盲评材料，但尚未根据自动指标推断作曲用途。已完成的实验支持关于符号控制与记谱导出的结论，艺术价值以及对独立乐谱的迁移能力仍待验证。
